% SOLO P2.
\section{Discusión y limitaciones}
\label{sec:discusion}

\subsection{Interpretación de los resultados}
A primera vista, la exactitud en prueba apenas varía entre el clasificador de referencia (0.735), la
regresión logística (0.743) y el árbol de decisión (0.744). Sin embargo, los dos modelos supervisados
sí aportan valor predictivo: pasan de detectar 0 abandonos a identificar entre el 72\% y el 78\% de
los 374 abandonos reales (292 y 270 clientes, respectivamente), cambiando aciertos sobre quienes
permanecen por detección temprana de riesgo. La regresión logística supera al árbol en \emph{Recall}
(+0.059 [0.016, 0.103]) y ROC-AUC (+0.023 [0.009, 0.040]), mientras que en F1 (+0.018 [$-$0.009,
0.049]), \emph{Precision} y exactitud las diferencias caen dentro del ruido muestral. Tres razones
explican este comportamiento:

\begin{enumerate}
    \item \emph{Señal mayormente aditiva y monótona.} El EDA mostró que las diferencias por contrato
    se mantienen dentro de cada tramo de antigüedad y que el abandono decrece sostenidamente con los
    meses de permanencia. Un modelo lineal en los \emph{log-odds} representa esa estructura con 46
    parámetros globales estimados sobre los 4225 clientes de entrenamiento.

    \item \emph{Fragmentación de datos en el árbol.} Cada división reduce el tamaño muestral en las
    ramas hasta hojas con un mínimo de 20 clientes. En prueba, el árbol genera alrededor de 100
    probabilidades distintas frente a 1407 de la regresión logística, lo que produce un ordenamiento
    más grueso (menor ROC-AUC) y una mayor brecha de F1 entre entrenamiento y validación (0.066
    frente a 0.012). De hecho, un árbol compacto de 8 hojas (\var{max\_depth=3}) alcanza F1 = 0.608 en
    validación, estadísticamente indistinguible del árbol elegido de 132 hojas (diferencia 0.011,
    IC 95\% [$-$0.014, 0.035]).

    \item \emph{Aporte limitado de interacciones.} La ventaja teórica del árbol solo se aprecia en
    nichos específicos: entre clientes sin internet detecta 11 de 25 abandonos, frente a ninguno en
    la regresión logística, pero son pocos casos para alterar el balance global.
\end{enumerate}

Además, la diferencia entre familias es modesta porque el techo lo impone la información disponible.
Ambos modelos coinciden en el 84.7\% de las predicciones de prueba y fallan conjuntamente en 61
abandonos (16.3\%): estos clientes presentan una antigüedad mediana de 35 meses, 44.3\% de contrato
mensual y 34.4\% de fibra óptica, un perfil casi idéntico al de quienes permanecen (38 meses,
42.9\% y 34.9\%). Distinguirlos requeriría variables ausentes en el conjunto, como reclamos, fallas
técnicas u ofertas de la competencia.

Finalmente, el tratamiento del desbalance (\var{class\_weight=\textquotedbl{}balanced\textquotedbl{}},
que penaliza 2.77 veces más los falsos negativos) tuvo un impacto mucho mayor que la fuerza de
regularización $C$. Permitió elevar sustancialmente el \emph{Recall}, pero llevó a señalar al 40.6\%
de los clientes de prueba (572 alertas, la mitad falsas alarmas) y descalibró las probabilidades
estimadas (media de 0.414 en regresión logística y 0.389 en el árbol frente a la tasa real de
0.265).

\subsection{Limitaciones}
\textbf{Del conjunto de datos.} Es una muestra de ejemplo de una empresa ficticia~\cite{ibm} sin
especificación del muestreo ni fechas de registro. Al ser una fotografía transversal, no permite
realizar una validación temporal (entrenar en el pasado y evaluar en el futuro) ni descartar por
completo una fuga temporal de origen; además, los resultados no son generalizables a un operador
real ni al mercado peruano. La antigüedad (\var{tenure}) padece sesgo de supervivencia: describe a
quienes no abandonaron antes, no la evolución causal de un cliente nuevo. Asimismo, la imputación con
cero de los 11 valores faltantes de \var{TotalCharges} (\var{tenure=0}) constituye un supuesto
metodológico no verificable.

\textbf{De la metodología.} Se empleó una sola partición estratificada (\emph{hold-out}, semilla 42).
Con 1409 observaciones en validación, diferencias de 0.002 a 0.01 en F1 quedan dentro del ruido
muestral y la configuración óptima podría cambiar con otra semilla; además, elegir el máximo entre 36
o 42 configuraciones introduce un ligero optimismo en validación respecto de prueba (0.625 a 0.617 en
regresión logística; 0.619 a 0.599 en árbol). El umbral de decisión se mantuvo fijo en 0.5 al no
contar con una matriz de costos reales de retención frente a pérdida de clientes, y solo se evaluaron
dos familias de modelos sin ensambles ni ingeniería de interacciones. Por último, la codificación
\emph{one-hot} completa introduce columnas redundantes (impidiendo interpretar coeficientes
individuales sueltos, razón por la que se reportan contrastes de \emph{odds}) y la correlación entre
\var{tenure} y \var{TotalCharges} (0.834 en entrenamiento) puede subestimar la importancia por
permutación de cada variable por separado.

\textbf{De los resultados.} Los intervalos de confianza en prueba son amplios ($\approx\pm$0.04 en F1)
y varios subgrupos tienen muy pocos abandonos (9 con contrato de dos años, 25 sin internet), lo que
hace imprecisas sus tasas. Las razones de \emph{odds} y reglas reflejan asociaciones observacionales,
no causalidad, y las probabilidades generadas sirven para ordenar por riesgo pero no están calibradas.

\subsection{Sesgos, ética y generalización}
El análisis por subgrupos muestra que \var{gender} es prácticamente neutro (OR = 1.05; \emph{Recall}
de 0.782 frente a 0.779 y tasa de falsos positivos de 0.271 frente a 0.270 en regresión logística).
En cambio, en \var{SeniorCitizen=1} (adultos mayores) la tasa de falsos positivos casi se duplica
respecto del resto (0.476 frente a 0.243 en regresión logística; 0.411 frente a 0.226 en árbol), a la
vez que suben el \emph{Recall} (0.929 frente a 0.728) y la \emph{Precision} (0.607 frente a 0.476)
debido a su mayor tasa base de abandono (44.1\% frente a 23.3\% en prueba).

El efecto directo de la variable es pequeño (OR = 1.18): la disparidad surge sobre todo porque, en el
conjunto completo, los adultos mayores presentan con mucha mayor frecuencia contrato mensual (70.7\%
frente a 52.0\%) y fibra óptica (72.8\% frente a 38.4\%). Por tanto, excluir \var{SeniorCitizen} del
entrenamiento no eliminaría la brecha, pues el contrato y el tipo de internet actuarían como
variables sustitutas.

Las implicancias éticas dependen de cómo se operacionalice la alerta. Si activa beneficios o
descuentos de retención, los adultos mayores resultan favorecidos y quedarían desatendidos quienes
tienen contratos largos (ningún modelo detectó los 9 abandonos con contrato de dos años). Si, por el
contrario, la alerta motivara restricciones o tratos desfavorables, perjudicaría de forma
desproporcionada a los adultos mayores. Asimismo, dado que el protocolo incluye atributos personales
(\var{gender}, \var{SeniorCitizen}, \var{Partner} y \var{Dependents}), un despliegue real exigiría
auditar su pertinencia normativa, validar con datos temporales propios del operador y monitorear
continuamente los errores por grupo.
